\subsection{Commutation recovered from a regularized filter}
\label{sec:shifted_density_commutation}

\begin{theorem}[Commutants of nonzero matrix powers]
    \label{thm:shifted_density_commutation}
    \lean{Matrix.PosDef.commute_rpow_iff,
        Matrix.PosSemidef.commute_add_smul_one_rpow_iff,
        Matrix.PosSemidef.commute_of_commute_shifted_inverse_power}
    \leanok
    Let $A$ be a positive definite complex matrix, $r\in\R\setminus\{0\}$,
    and $\rho$ any complex matrix of the same size. Then $A^r$ commutes with
    $\rho$ if and only if $A$ commutes with $\rho$.
    If $x\ge0$ and $b>0$, then $(x+b\Id)^r$ commutes with $\rho$ if and only
    if $x$ commutes with $\rho$. Consequently, for $a>0$, commutation of $\rho$ with the
    regularized filter $(x+b\Id)^{-a/2}$ implies $[\rho,x]=0$.
    This is the functional-calculus implication in
    \cite[Section~4, descending stationarity]{OpenAI2026PolynomialPEPS}.
\end{theorem}
\begin{proof}
    \leanok
    Functional calculus preserves commutation with an arbitrary matrix.
    The power identity $(A^r)^{1/r}=A$ gives the converse.
    For the shifted matrix, subtracting the scalar matrix $b\Id$ preserves
    commutation, and $-a/2\ne0$. Taking $r=2$ also recovers commutation
    with a positive definite filter $L$ from commutation with $L^2$.
\end{proof}
